\section*{الفصل \ref{ch:g12:destruction-of-mountains} --- هدم الجبال}

\begin{solution}{exo:g12:destruction-of-mountains:1}
(أ) \omterm{def:g12:destruction-of-mountains:denudation}{معدل التجريد}؛ (ب) \omterm{def:g12:destruction-of-mountains:load}{الحمولة القاعية}؛ (ج) \omterm{def:g12:destruction-of-mountains:load}{الحمولة الذائبة}؛
(د) \omterm{def:g12:destruction-of-mountains:diagenesis}{التصخّر}؛ (ه) \omterm{def:g12:destruction-of-mountains:rebound}{الارتداد الأيزوستاتي}.
\end{solution}

\begin{solution}{exo:g12:destruction-of-mountains:2}
(1) رصيص. (2) حجر رملي. (3) سجيل. (4) بريشيا.
\end{solution}

\begin{solution}{exo:g12:destruction-of-mountains:3}
\qty{1}{m/Myr} تساوي \qty{1000}{mm} في \num{1000000} سنة، أي \qty{0.001}{mm} في
السنة. (1) \qty{0.054}{mm} في السنة. (2) \qty{0.218}{mm} في السنة.
(3) $\qty{10}{mm}$ في السنة تساوي \qty{10}{m} في 1000 سنة، أي \qty{10000}{m/Myr}.
\end{solution}

\begin{solution}{exo:g12:destruction-of-mountains:4}
الحمولة الصلبة \qty{85}{\percent} من المجموع، فالمجموع $60/0.85
\approx 70.6$ مليون طن، \omterm{def:g12:destruction-of-mountains:load}{والحمولة الذائبة} نحو $70.6 - 60 \approx
10.6$ مليون طن في السنة.
\end{solution}

\begin{solution}{exo:g12:destruction-of-mountains:5}
تجوية الغرانيت، فالنقل بنهر، فالترسيب في حوض،
فالدفن تحت طبقات أحدث، فالالتحام.
\end{solution}

\begin{solution}{exo:g12:destruction-of-mountains:6}
عند \qty{30}{cm/s} تُرفع حبيبة الرمل (فالسرعة فوق \qty{20}{cm/s})؛ وتبقى
حبيبة الغرين عالقة. وعند \qty{5}{cm/s} تترسب حبيبة الرمل
(فالسرعة دون \qty{10}{cm/s})، لكن حبيبة الغرين تظل عالقة، لأن
\qty{0.1}{cm/s} تكفي لإبقائها كذلك.
\end{solution}

\begin{solution}{exo:g12:destruction-of-mountains:7}
الحمولة الكلية $4/0.85 \approx 4.71$ مليون طن، أي $\num{4.71e9}/2830 \approx
\qty{1.66e6}{m^3}$ من الصخر في السنة. وعلى $\qty{5e10}{m^2}$: $\num{3.3e-5}$~\unit{m} في
السنة، أي نحو \qty{33}{m/Myr}، وهذا دون الوسيط البالغ \qty{54}{m/Myr}.
\end{solution}

\begin{solution}{exo:g12:destruction-of-mountains:8}
(1) $1 - 2700/3300 \approx 0.182$. (2) نحو \qty{182}{m}. (3) أكثر قليلًا
من \qty{163}{m} التي في الفصل: فكلما كانت القشرة أخف بالنسبة إلى
الوشاح، قلّ ارتفاع العمود.
\end{solution}

\begin{solution}{exo:g12:destruction-of-mountains:9}
(1) \qty{8}{km} من الصخر. (2) $8000 \times 0.163 \approx \qty{1300}{m}$.
(3) $8000/200 = 40$ مليون سنة.
\end{solution}

\begin{solution}{exo:g12:destruction-of-mountains:10}
(1) $794/43 \approx 18$ مليون سنة. (2) لا ينخفض السطح إلا بمقدار
$0.163$ من السمك المُعرّى، فالصخر الواجب إزالته $794/0.163
\approx \qty{4900}{m}$، وهذا يستغرق نحو $4900/43 \approx 110$ ملايين سنة.
(3) القارات لا تُبرى فحسب: فبناء الجبال وصهارية
الأقواس يظلان يرفعانها ويضيفان إليها قشرة جديدة.
\end{solution}

\begin{solution}{exo:g12:destruction-of-mountains:11}
(1) $\num{1e14}/2830 \approx \qty{3.5e10}{m^3}$، أي نحو
\qty{35}{km^3}. (2) $100/15.5 \approx 6.5$ سنة.
\end{solution}

\begin{solution}{exo:g12:destruction-of-mountains:12}
$E = kh$ و$h = h_0 e^{-fkt}$ تعطيان $E(t) = kh_0 e^{-fkt}$. والسمك
المُزال من $t = 0$ إلى $t \to \infty$ هو $\int_0^{\infty} kh_0
e^{-fkt}\,\mathrm{d}t = kh_0/(fk) = h_0/f$. فلخفض السطح بمقدار $h_0$،
يجب أن تزيل التعرية $h_0/f \approx 6.1\,h_0$، كما تقول قضية
الارتداد: فكل متر يُعرّى لا يخفض السطح إلا بمقدار $f$ متر.
\end{solution}

\begin{solution}{exo:g12:destruction-of-mountains:13}
(1) $10 \times 1000 = \qty{10000}{mm}$، أي \qty{10}{m}. (2) إزالة حمل من
جليد سمكه $t$ تكافئ إزالة كتلة $\rho_{\text{جليد}}\,t$ لكل وحدة
مساحة؛ فيرتفع العمود إلى أن يزيح من الوشاح أقل بهذا القدر نفسه:
$2000 \times 900/3380 \approx \qty{530}{m}$. (3) ترتفع الأرض أقل فأقل
كلما اقتربت من توازنها الجديد: فما يحرك الارتفاع هو الجزء الذي لم
يُستعد بعد.
\end{solution}

\begin{solution}{exo:g12:destruction-of-mountains:14}
(1) قبل 1900: نحو \qty{56}{m/Myr} (كما في المثال). وفي
1987--2006: $145/0.85 \approx 171$ مليون طن، و$\num{1.71e11}/2830
\approx \qty{6.0e7}{m^3}$، على $\qty{2.98e12}{m^2}$: أي نحو
\qty{20}{m/Myr}. (2) القيمة الأقدم: فالحمولة اللاحقة تحجزها السدود
ومنشآت أخرى. (3) لا: فالنقص $400 - 145 = 255$ مليون طن في
السنة، لا تفسر السدود منه إلا 100 إلى 150، أي نحو النصف.
\end{solution}

\begin{solution}{exo:g12:destruction-of-mountains:15}
(1) \qty{6}{km} على الأقل. (2) $310 - 252 = 58$ مليون سنة على الأكثر، إذن
$6000/58 \approx \qty{100}{m/Myr}$ على الأقل. (3) يقع بين الوسيط
(\qty{54}{m/Myr}) والمتوسط (\qty{218}{m/Myr}) للأحواض: إنه معدل
عادي لحزام جبلي.
\end{solution}

\begin{solution}{pb:g12:destruction-of-mountains:1}
\textbf{1.} $E = 0.209 \times 2000 \approx \qty{418}{m/Myr}$.

\textbf{2.} \qty{0.418}{mm} في السنة.

\textbf{3.} أسرع من كليهما: إنها سلسلة فتية شديدة الانحدار، سريعة سرعة
الأحواض النشطة زلزاليًا وأكثر.

\textbf{4.} متوسط انحداره.

\textbf{5.} طبقة سمكها $\qty{0.418e-3}{m}$ على $\qty{1e6}{m^2}$ تساوي
\qty{418}{m^3}، و$418 \times 2830 \approx \qty{1.18e6}{kg}$: أي نحو 1200
طن.

\textbf{6.} $2000/418 \approx 4.8$ مليون سنة.

\textbf{7.} $h(t) = 2000\,e^{-0.209\,t}$ (حيث $t$ بملايين السنين).

\textbf{8.} $1/0.209 \approx 4.8$ مليون سنة.

\textbf{9.} $\ln 2/0.209 \approx 3.3$ مليون سنة.

\textbf{10.} $e^{-0.209 \times 11} = e^{-2.30} \approx 0.100$.

\textbf{11.} $1 - 2830/3380 \approx 0.163$.

\textbf{12.} ينخفض السطح بمقدار $100 \times 0.163 \approx \qty{16}{m}$؛
ويرتفع \omterm{def:g12:continental-crust:root}{الجذر} بنحو \qty{84}{m}.

\textbf{13.} $h' = -0.163 \times 0.209\,h \approx -0.0341\,h$، إذن $h(t) =
2000\,e^{-0.0341\,t}$.

\textbf{14.} $(2000 - 200)/0.163 \approx \qty{11000}{m}$ من الصخر.

\textbf{15.} صخور كانت على عمق نحو \qty{11}{km}: صخور متحولة
مثل الشيست والنيس، وغرانيتات.

\textbf{16.} $1/0.0341 \approx 29$ مليون سنة.

\textbf{17.} $29 \times \ln 10 \approx 67$ مليون سنة.

\textbf{18.} $2000\,e^{-0.0341 \times 200} \approx 2000 \times 0.0011 \approx
\qty{2}{m}$.

\textbf{19.} $2000\,e^{-0.0341 \times 250} \approx 2000 \times 0.0002 \approx
\qty{0.4}{m}$: سهل مستوٍ. وقمة ارتفاعها \qty{2037}{m} تعني أن النموذج
وحده لا يطابق الواقع: لا بد أن المنطقة رُفعت من جديد منذ ذلك الحين، أو تعرّت
أبطأ مما يفترض النموذج.

\textbf{20.} ينتصف متوسط الارتفاع في $29 \times \ln 2 \approx$
\textbf{20 مليون سنة}.
\end{solution}
